% chapters/99-vita-1page.tex
% -----------------------------------------------------------------------------
% VITA (condensed one-page narrative version)
% -----------------------------------------------------------------------------
\chapter*{VITA}
\addcontentsline{toc}{chapter}{VITA}
\markboth{VITA}{VITA}

Dan Mailman is a PhD candidate in Computer Science at the University of Tennessee at Chattanooga. His doctoral work focuses on gesture-efficient global language input methods: how people can enter Unicode-rich text---including Mandarin, IPA, and other multilingual writing systems---quickly and accurately on modern touch devices. His dissertation, \emph{Gesture-Efficient Global Language Input Methods}, develops a practical engineering and evaluation framework for high-efficiency input, emphasizing measurable productivity, learnability for touch typing, and reduced gesture effort.



Alongside the research, he builds end-to-end prototypes that connect theory to usable software. These projects include a family of grid-based input interfaces (including the \emph{\textmu Lex} framework and the \emph{muCiYu} Mandarin prototype), with careful attention to interaction design, error handling, and reproducible evaluation. He also maintains a structured writing and tooling pipeline for the dissertation using LaTeX, XeLaTeX, and Quarto, and routinely develops supporting analyses and visualizations in R and Python.



Professionally, Dan works in education technology and has been affiliated with Great Minds PBC while completing the degree. He is also exploring commercialization paths for his input-method research through Language Genies, Inc., and he publishes technical and creative work under the Studio MindStride umbrella. Outside of his research and writing, he enjoys cycling, rowing, and hosting film discussions.

